\chapter{网格是几何与离散算子的接口}\label{ch:20}
\begin{chaptersummary}
模型和离散空间确定后，网格将物理几何、边界和局部尺度引入算子。本章从参考单元映射出发，解释形状与尺寸如何改变梯度、积分及界面配对，再联系曲边几何和守恒检验。这些几何相容性要求也构成下一章复杂网格技术的共同基础。
\end{chaptersummary}
\learninggoals{计算参考到物理单元的梯度与积分变换；区分几何误差与逼近误差；说明网格形状、邻接及边界表示对空间算子的影响。}
\studyroute{略去高维网格生成、几何近似误差证明和矩阵实现；本章保留算子结构，代数求解与并行数据布局见第 0 卷。}{element mapping; mesh quality; curved element; metric identity; discrete geometric conservation}{比较等边三角形与细长三角形的 Jacobian 和线性有限元刚度矩阵，解释梯度放大来自哪一个几何因子。}
\section{几何、邻接与离散耦合}\index{网格几何}
网格决定几何表示误差、局部离散误差、稳定性限制和矩阵非零结构。几何处理与网格质量控制是计算工作流的重要环节。

\begin{authornote}
\textbf{潜在疑问：}
“为什么网格不仅影响几何外形，还会影响矩阵结构、稳定性限制和求解器行为？”

\textbf{作者补充：}
网格不仅给出采样位置，还规定单元之间如何交换信息。邻接关系决定离散耦合，单元形状和尺寸进入导数、积分与时间步长估计，边界表示影响边界条件施加。同一连续方程在不同网格上会产生不同的离散耦合和误差。

\textbf{核心概念联系：}
\[
\text{mesh geometry}
\longrightarrow
\text{discrete coupling}
\longrightarrow
\text{matrix and solver behavior}
\]
\end{authornote}

贴体网格使网格边界与物体边界一致。其边界条件施加自然，适合解析壁面边界层，但复杂几何和大运动会导致网格生成与质量控制困难。

\section{参考单元映射与几何尺度}
设 $x=F_K(\xi)$，$A_K=\partial F_K/\partial\xi$ 可逆，则
\[
\nabla_xu=A_K^{-\mathsf T}\nabla_\xi\widehat u,\qquad
\dd x=|\det A_K|\dd\xi.
\]
Poisson 刚度条目成为
\[
K_{ij}^{(K)}
=\int_{\widehat K}(\nabla_\xi\phi_i)^{\mathsf T}
A_K^{-1}A_K^{-\mathsf T}\nabla_\xi\phi_j\,|\det A_K|\dd\xi.
\]
细长或退化单元改变映射奇异值，放大某些方向的梯度及误差常数。
结构化网格具有规则连接；非结构网格允许更灵活的局部形状与邻接。
两者都需规定边界取向、面测度和离散积分。

\section{曲边单元与几何误差}
高阶场近似常需相应阶数的曲边几何。
参考到物理映射的次数、体积求积和面求积共同影响总误差。
在守恒律中，常数解保持要求离散几何因子满足对应的度量恒等式；
共享面几何和相反法向是通量抵消的基础。
验证时可先用常数场、低阶多项式和解析曲面检查几何积分，再评估物理解的收敛。

\aiexercise{对映射 $x=A\xi+b$，分别取旋转、各向同性缩放和强拉伸的 $A$。推导其梯度、质量与刚度缩放，并比较三种网格上的同一线性场误差。}
{将一个圆边界分别用折线和二次曲线表示，固定场近似阶数，设计区分几何误差与空间逼近误差的收敛实验。}
